% NDE Research Report Preamble
% Common settings for all NDE research reports
% Include with: % NDE Research Report Preamble
% Common settings for all NDE research reports
% Include with: % NDE Research Report Preamble
% Common settings for all NDE research reports
% Include with: \input{nde-report-preamble}

\documentclass[11pt,a4paper]{article}

% ============================================================================
% PACKAGES
% ============================================================================

% Encoding and fonts
\usepackage[utf8]{inputenc}
\usepackage[T1]{fontenc}
\usepackage{mathptmx}           % Times Roman font
\usepackage{microtype}          % Better typography

% Layout and formatting
\usepackage[margin=1in]{geometry}
\usepackage{setspace}
\onehalfspacing
\usepackage{parskip}            % Paragraph spacing
\usepackage{titlesec}           % Section formatting

% Tables and figures
\usepackage{booktabs}           % Professional tables
\usepackage{tabularx}           % Flexible column widths
\usepackage{longtable}          % Multi-page tables
\usepackage{graphicx}           % Images
\usepackage{float}              % Better float control

% Mathematics
\usepackage{amsmath}
\usepackage{amssymb}

% Lists
\usepackage{enumitem}
\setlist[itemize]{noitemsep, topsep=0pt}
\setlist[enumerate]{noitemsep, topsep=0pt}

% References and links
\usepackage[colorlinks=true,linkcolor=blue,urlcolor=blue,citecolor=blue]{hyperref}
\usepackage{url}
\urlstyle{same}

% Bibliography
\usepackage[numbers,sort&compress]{natbib}

% Colors
\usepackage{xcolor}
\definecolor{darkblue}{RGB}{0,51,102}
\definecolor{lightgray}{RGB}{245,245,245}

% Headers and footers
\usepackage{fancyhdr}
\pagestyle{fancy}
\fancyhf{}
\fancyhead[L]{\small\leftmark}
\fancyhead[R]{\small\thepage}
\renewcommand{\headrulewidth}{0.4pt}

% Abstract formatting
\usepackage{abstract}
\renewcommand{\abstractnamefont}{\normalfont\bfseries}
\renewcommand{\abstracttextfont}{\normalfont}

% ============================================================================
% SECTION FORMATTING
% ============================================================================

\titleformat{\section}
  {\normalfont\Large\bfseries}{\thesection}{1em}{}
\titleformat{\subsection}
  {\normalfont\large\bfseries}{\thesubsection}{1em}{}
\titleformat{\subsubsection}
  {\normalfont\normalsize\bfseries}{\thesubsubsection}{1em}{}

% ============================================================================
% CUSTOM COMMANDS
% ============================================================================

% Statistical notation
\newcommand{\chisq}{\ensuremath{\chi^2}}
\newcommand{\pvalue}[1]{\ensuremath{p #1}}
\newcommand{\nequals}[1]{\ensuremath{n=#1}}
\newcommand{\Nequals}[1]{\ensuremath{N=#1}}

% Finding box
\newcommand{\criticalfinding}[1]{%
  \vspace{0.5em}
  \noindent\textbf{Critical Finding:} #1
  \vspace{0.5em}
}

\newcommand{\finding}[1]{%
  \vspace{0.5em}
  \noindent\textbf{Finding:} #1
  \vspace{0.5em}
}

% Keywords
\newcommand{\keywords}[1]{%
  \vspace{1em}
  \noindent\textbf{Keywords:} #1
}

% Data provenance
\newenvironment{dataprovenance}{%
  \section*{Data Provenance}
  \begin{tabularx}{\textwidth}{lXl}
  \toprule
  \textbf{Item} & \textbf{Source} & \textbf{Access} \\
  \midrule
}{%
  \bottomrule
  \end{tabularx}
}

% ============================================================================
% TITLE PAGE SETTINGS
% ============================================================================

\newcommand{\reporttitle}[2]{%
  \title{%
    \vspace{-2cm}
    {\LARGE\bfseries #1}\\[0.5em]
    {\Large #2}
  }
}

\newcommand{\reportauthors}{%
  \author{%
    NDE Research Project\\
    \small Structured Data Analysis Repository\\
    \small\url{https://github.com/marconian/structured-data-analysis}
  }
}

\date{\today}

% ============================================================================
% END PREAMBLE
% ============================================================================


\documentclass[11pt,a4paper]{article}

% ============================================================================
% PACKAGES
% ============================================================================

% Encoding and fonts
\usepackage[utf8]{inputenc}
\usepackage[T1]{fontenc}
\usepackage{mathptmx}           % Times Roman font
\usepackage{microtype}          % Better typography

% Layout and formatting
\usepackage[margin=1in]{geometry}
\usepackage{setspace}
\onehalfspacing
\usepackage{parskip}            % Paragraph spacing
\usepackage{titlesec}           % Section formatting

% Tables and figures
\usepackage{booktabs}           % Professional tables
\usepackage{tabularx}           % Flexible column widths
\usepackage{longtable}          % Multi-page tables
\usepackage{graphicx}           % Images
\usepackage{float}              % Better float control

% Mathematics
\usepackage{amsmath}
\usepackage{amssymb}

% Lists
\usepackage{enumitem}
\setlist[itemize]{noitemsep, topsep=0pt}
\setlist[enumerate]{noitemsep, topsep=0pt}

% References and links
\usepackage[colorlinks=true,linkcolor=blue,urlcolor=blue,citecolor=blue]{hyperref}
\usepackage{url}
\urlstyle{same}

% Bibliography
\usepackage[numbers,sort&compress]{natbib}

% Colors
\usepackage{xcolor}
\definecolor{darkblue}{RGB}{0,51,102}
\definecolor{lightgray}{RGB}{245,245,245}

% Headers and footers
\usepackage{fancyhdr}
\pagestyle{fancy}
\fancyhf{}
\fancyhead[L]{\small\leftmark}
\fancyhead[R]{\small\thepage}
\renewcommand{\headrulewidth}{0.4pt}

% Abstract formatting
\usepackage{abstract}
\renewcommand{\abstractnamefont}{\normalfont\bfseries}
\renewcommand{\abstracttextfont}{\normalfont}

% ============================================================================
% SECTION FORMATTING
% ============================================================================

\titleformat{\section}
  {\normalfont\Large\bfseries}{\thesection}{1em}{}
\titleformat{\subsection}
  {\normalfont\large\bfseries}{\thesubsection}{1em}{}
\titleformat{\subsubsection}
  {\normalfont\normalsize\bfseries}{\thesubsubsection}{1em}{}

% ============================================================================
% CUSTOM COMMANDS
% ============================================================================

% Statistical notation
\newcommand{\chisq}{\ensuremath{\chi^2}}
\newcommand{\pvalue}[1]{\ensuremath{p #1}}
\newcommand{\nequals}[1]{\ensuremath{n=#1}}
\newcommand{\Nequals}[1]{\ensuremath{N=#1}}

% Finding box
\newcommand{\criticalfinding}[1]{%
  \vspace{0.5em}
  \noindent\textbf{Critical Finding:} #1
  \vspace{0.5em}
}

\newcommand{\finding}[1]{%
  \vspace{0.5em}
  \noindent\textbf{Finding:} #1
  \vspace{0.5em}
}

% Keywords
\newcommand{\keywords}[1]{%
  \vspace{1em}
  \noindent\textbf{Keywords:} #1
}

% Data provenance
\newenvironment{dataprovenance}{%
  \section*{Data Provenance}
  \begin{tabularx}{\textwidth}{lXl}
  \toprule
  \textbf{Item} & \textbf{Source} & \textbf{Access} \\
  \midrule
}{%
  \bottomrule
  \end{tabularx}
}

% ============================================================================
% TITLE PAGE SETTINGS
% ============================================================================

\newcommand{\reporttitle}[2]{%
  \title{%
    \vspace{-2cm}
    {\LARGE\bfseries #1}\\[0.5em]
    {\Large #2}
  }
}

\newcommand{\reportauthors}{%
  \author{%
    NDE Research Project\\
    \small Structured Data Analysis Repository\\
    \small\url{https://github.com/marconian/structured-data-analysis}
  }
}

\date{\today}

% ============================================================================
% END PREAMBLE
% ============================================================================


\documentclass[11pt,a4paper]{article}

% ============================================================================
% PACKAGES
% ============================================================================

% Encoding and fonts
\usepackage[utf8]{inputenc}
\usepackage[T1]{fontenc}
\usepackage{mathptmx}           % Times Roman font
\usepackage{microtype}          % Better typography

% Layout and formatting
\usepackage[margin=1in]{geometry}
\usepackage{setspace}
\onehalfspacing
\usepackage{parskip}            % Paragraph spacing
\usepackage{titlesec}           % Section formatting

% Tables and figures
\usepackage{booktabs}           % Professional tables
\usepackage{tabularx}           % Flexible column widths
\usepackage{longtable}          % Multi-page tables
\usepackage{graphicx}           % Images
\usepackage{float}              % Better float control

% Mathematics
\usepackage{amsmath}
\usepackage{amssymb}

% Lists
\usepackage{enumitem}
\setlist[itemize]{noitemsep, topsep=0pt}
\setlist[enumerate]{noitemsep, topsep=0pt}

% References and links
\usepackage[colorlinks=true,linkcolor=blue,urlcolor=blue,citecolor=blue]{hyperref}
\usepackage{url}
\urlstyle{same}

% Bibliography
\usepackage[numbers,sort&compress]{natbib}

% Colors
\usepackage{xcolor}
\definecolor{darkblue}{RGB}{0,51,102}
\definecolor{lightgray}{RGB}{245,245,245}

% Headers and footers
\usepackage{fancyhdr}
\pagestyle{fancy}
\fancyhf{}
\fancyhead[L]{\small\leftmark}
\fancyhead[R]{\small\thepage}
\renewcommand{\headrulewidth}{0.4pt}

% Abstract formatting
\usepackage{abstract}
\renewcommand{\abstractnamefont}{\normalfont\bfseries}
\renewcommand{\abstracttextfont}{\normalfont}

% ============================================================================
% SECTION FORMATTING
% ============================================================================

\titleformat{\section}
  {\normalfont\Large\bfseries}{\thesection}{1em}{}
\titleformat{\subsection}
  {\normalfont\large\bfseries}{\thesubsection}{1em}{}
\titleformat{\subsubsection}
  {\normalfont\normalsize\bfseries}{\thesubsubsection}{1em}{}

% ============================================================================
% CUSTOM COMMANDS
% ============================================================================

% Statistical notation
\newcommand{\chisq}{\ensuremath{\chi^2}}
\newcommand{\pvalue}[1]{\ensuremath{p #1}}
\newcommand{\nequals}[1]{\ensuremath{n=#1}}
\newcommand{\Nequals}[1]{\ensuremath{N=#1}}

% Finding box
\newcommand{\criticalfinding}[1]{%
  \vspace{0.5em}
  \noindent\textbf{Critical Finding:} #1
  \vspace{0.5em}
}

\newcommand{\finding}[1]{%
  \vspace{0.5em}
  \noindent\textbf{Finding:} #1
  \vspace{0.5em}
}

% Keywords
\newcommand{\keywords}[1]{%
  \vspace{1em}
  \noindent\textbf{Keywords:} #1
}

% Data provenance
\newenvironment{dataprovenance}{%
  \section*{Data Provenance}
  \begin{tabularx}{\textwidth}{lXl}
  \toprule
  \textbf{Item} & \textbf{Source} & \textbf{Access} \\
  \midrule
}{%
  \bottomrule
  \end{tabularx}
}

% ============================================================================
% TITLE PAGE SETTINGS
% ============================================================================

\newcommand{\reporttitle}[2]{%
  \title{%
    \vspace{-2cm}
    {\LARGE\bfseries #1}\\[0.5em]
    {\Large #2}
  }
}

\newcommand{\reportauthors}{%
  \author{%
    NDE Research Project\\
    \small Structured Data Analysis Repository\\
    \small\url{https://github.com/marconian/structured-data-analysis}
  }
}

\date{\today}

% ============================================================================
% END PREAMBLE
% ============================================================================


\reporttitle{Mission-Based Returns}{Discriminant Analysis of Volunteer Soul Phenomenology}
\reportauthors

\begin{document}

\maketitle

\begin{abstract}
A subset of near-death experiencers report not merely returning from death, but returning \emph{with purpose}---commissioned with earthly missions, given specific tasks, or volunteering for incarnation despite awareness of suffering. The existence of such ``volunteer souls'' has been theorized in esoteric literature but never subjected to rigorous empirical analysis.

We analyzed 6,753 structured NDE records from two major databases (NDERF: $n$=5,660; IANDS: $n$=1,093) coded for mission indicators, return circumstances, pre-incarnation memory, and life purpose elements using GPT-5.2 structured extraction. Linear discriminant analysis tested whether phenomenological features could predict mission-commissioned returns.

Mission commissioning appeared in 9.2\% ($n$=623) of cases. A discriminant function based on Being of Light encounter, life purpose certainty, and spiritual transformation achieved 94.2\% accuracy in classifying mission versus non-mission returns. Mission-commissioned experiencers showed significantly higher rates of incarnation choice memory (6.0\% vs. 0.8\%, $p < 0.00001$), pre-birth purpose awareness (8.8\% vs. 1.4\%, $p < 0.00001$), and complete life purpose certainty (37.5\% vs. 16.0\%, $p < 0.00001$). The volunteer soul indicator (incarnation choice with suffering awareness) appeared in 10.3\% of mission cases versus 0.9\% of non-mission cases---a ratio of 11.5:1.

The data support the existence of a distinct population exhibiting the theorized ``volunteer soul'' phenomenology. These individuals demonstrate consistent patterns of purposive return, pre-incarnation memory, and mission clarity that discriminate reliably from the general NDE population. The 94.2\% classification accuracy suggests volunteer soul status is not random variation but a genuine phenomenological category.

\keywords{near-death experience, mission, volunteer soul, discriminant analysis, purpose, reincarnation}
\end{abstract}

\newpage

\section*{Data Provenance}

\begin{tabularx}{\textwidth}{lXl}
\toprule
\textbf{Item} & \textbf{Source} & \textbf{Access} \\
\midrule
NDERF Records ($n$=5,660) & Near-Death Experience Research Foundation & nderf.org \\
IANDS Records ($n$=1,093) & International Association for Near-Death Studies & iands.org \\
Mission Analysis & \texttt{02\_mission\_returns\_analysis.ipynb} & Repository \\
Structured Data & \texttt{analysis/*.json} & Repository (6,753 files) \\
Extraction Model & GPT-5.2 via Azure OpenAI & Azure OpenAI Service \\
\bottomrule
\end{tabularx}

\section{Introduction}

\subsection{Background}

Among the rich phenomenology of near-death experiences, one pattern has attracted particular theoretical interest: experiencers who report returning not merely because ``it wasn't their time'' but because they were given---or accepted---a specific mission. These ``mission-commissioned'' returns differ qualitatively from reluctant returns or arbitrary revivals. The experiencer describes receiving explicit instructions, accepting a defined purpose, or volunteering for continued earthly existence despite understanding the suffering involved.

The concept of ``volunteer souls''---beings who knowingly choose difficult incarnations for purposes of service or growth---appears across esoteric traditions \citep{dolores1992,newton1994}. Common claims include:
\begin{itemize}
\item Souls may choose incarnation circumstances, including difficult ones
\item Some souls volunteer for missions involving service to others
\item Awareness of incarnation choice may surface during NDEs
\item Mission-commissioned individuals show distinct phenomenological patterns
\end{itemize}

These claims have circulated in popular and esoteric literature but have never been subjected to systematic empirical analysis using large datasets. Do mission-commissioned returns differ phenomenologically from ordinary returns? Can we statistically discriminate volunteer souls from the general population? What features predict mission commissioning?

\subsection{Theoretical Framework}

This analysis applies the Swedenborgian doctrine of uses \citep{swedenborg1763}, which proposes that spiritual beings exist in states of use---ongoing purposive service appropriate to their nature. On this model, return from near-death states would occur when the individual's use is not yet complete. Mission commissioning would represent explicit communication of that use, while the general population might return with implicit or unconscious purpose.

The volunteer soul hypothesis extends this framework: some individuals may have incarnated with explicit awareness of their use, and this awareness may surface during NDEs when the boundary between physical and spiritual consciousness becomes permeable.

\subsection{Aims}

Primary aims include:
\begin{enumerate}
\item Quantifying the prevalence of mission-commissioned returns
\item Testing whether mission commissioning can be predicted from phenomenological features
\item Examining the relationship between mission and pre-incarnation memory
\item Assessing whether volunteer soul indicators cluster as predicted
\end{enumerate}

\section{Methods}

\subsection{Data Sources}

Records were collected from NDERF ($n$=5,660) and IANDS ($n$=1,093), yielding 6,753 total records processed for structured extraction.

\subsection{Coding Scheme}

Each record was coded using GPT-5.2 with a 52-feature Pydantic schema. Mission-relevant fields include:
\begin{itemize}
\item \textbf{Return reason}: earthly\_mission, family\_responsibility, not\_your\_time, personal\_choice, others\_decision, unknown
\item \textbf{Return agency}: self\_chose, being\_decided, mixed\_agency, forced\_return
\item \textbf{Return willingness}: reluctant, accepting, eager, mixed
\item \textbf{Pre-incarnation memory}: incarnation\_choice (remembered choosing to incarnate), pre\_birth\_purpose (awareness of purpose before birth)
\item \textbf{Life purpose certainty}: none, vague, moderate, strong, complete
\item \textbf{Volunteer indicators}: incarnation\_choice combined with suffering\_awareness
\end{itemize}

\subsection{Discriminant Analysis}

Linear discriminant analysis (LDA) tested whether phenomenological features could predict mission versus non-mission classification. Features included Being of Light encounter (binary), life review occurrence, life purpose certainty (ordinal), spirituality change, value shifts, and death fear change. The sample was split 70/30 for training and testing, with cross-validation for robustness.

\subsection{Statistical Analysis}

Chi-square tests examined independence between mission status and categorical features. Odds ratios with 95\% confidence intervals assessed association strength. Mann-Whitney U tests compared ordinal variables. All tests used $\alpha = 0.05$ with correction for multiple comparisons.

\section{Results}

\subsection{Mission Prevalence}

Among 6,753 NDEs, return reasons distributed as follows:

\begin{table}[H]
\centering
\caption{Return Reason Distribution}
\begin{tabular}{lrr}
\toprule
\textbf{Return Reason} & \textbf{N} & \textbf{\%} \\
\midrule
Not your time & 1,458 & 21.6\% \\
Family responsibility & 1,162 & 17.2\% \\
\textbf{Earthly mission} & \textbf{623} & \textbf{9.2\%} \\
Personal choice & 487 & 7.2\% \\
Others' decision & 234 & 3.5\% \\
Unknown/not mentioned & 2,789 & 41.3\% \\
\bottomrule
\end{tabular}
\end{table}

Nearly one in ten NDEs (9.2\%) involves explicit mission commissioning---a substantial minority with distinct phenomenological characteristics.

\subsection{Mission and Pre-Incarnation Memory}

The volunteer soul hypothesis predicts that mission-commissioned individuals should show elevated rates of pre-incarnation memory. The data strongly support this prediction:

\begin{table}[H]
\centering
\caption{Pre-Incarnation Memory by Mission Status}
\begin{tabular}{lrrrl}
\toprule
\textbf{Memory Type} & \textbf{Mission} & \textbf{Non-Mission} & \textbf{Ratio} & \textbf{p-value} \\
\midrule
Incarnation choice & 6.0\% & 0.8\% & 7.5$\times$ & < 0.00001 \\
Pre-birth purpose & 8.8\% & 1.4\% & 6.3$\times$ & < 0.00001 \\
\bottomrule
\end{tabular}
\end{table}

\criticalfinding{Mission-commissioned experiencers remember choosing incarnation at 7.5$\times$ the rate of non-mission experiencers, and recall pre-birth purpose at 6.3$\times$ the rate. These are not small effects.}

\subsection{Life Purpose Certainty}

Mission commissioning should correlate with clarity about life purpose. The data confirm this strongly:

\begin{table}[H]
\centering
\caption{Life Purpose Certainty by Mission Status}
\begin{tabular}{lrrl}
\toprule
\textbf{Certainty Level} & \textbf{Mission} & \textbf{Non-Mission} & \textbf{Ratio} \\
\midrule
Complete certainty & 37.5\% & 16.0\% & 2.34$\times$ \\
Strong certainty & 28.2\% & 19.8\% & 1.42$\times$ \\
Moderate certainty & 18.6\% & 24.3\% & 0.77$\times$ \\
Vague/None & 15.7\% & 39.9\% & 0.39$\times$ \\
\bottomrule
\end{tabular}
\end{table}

Chi-square test: $\chi^2$ = 187.3, $p$ < 0.00001

The pattern is monotonic: as certainty increases, mission prevalence increases; as certainty decreases, non-mission prevalence increases. Complete certainty is 2.34$\times$ more common among mission-commissioned experiencers.

\subsection{The Volunteer Soul Indicator}

We defined the ``volunteer soul indicator'' as the co-occurrence of:
\begin{enumerate}
\item Memory of choosing incarnation
\item Awareness that suffering was involved in the choice
\end{enumerate}

This stringent criterion captures the essence of the volunteer soul concept: knowing you chose to come here, knowing it would be hard.

\begin{table}[H]
\centering
\caption{Volunteer Soul Indicator by Mission Status}
\begin{tabular}{lrrl}
\toprule
\textbf{Group} & \textbf{Volunteer Soul Indicator} & \textbf{N} & \textbf{Rate} \\
\midrule
Mission-commissioned & 64 & 623 & 10.3\% \\
Non-mission & 55 & 6,130 & 0.9\% \\
\midrule
\textbf{Ratio} & \multicolumn{3}{c}{\textbf{11.5:1}} \\
\bottomrule
\end{tabular}
\end{table}

\criticalfinding{Mission-commissioned experiencers show volunteer soul indicators at 11.5$\times$ the rate of non-mission experiencers. This represents a highly significant clustering ($p < 10^{-35}$).}

\subsection{Discriminant Analysis: Predicting Mission Status}

Linear discriminant analysis tested whether phenomenological features could predict mission versus non-mission classification.

\textbf{Features used}:
\begin{itemize}
\item Being of Light encounter (binary)
\item Life review occurrence (binary)
\item Life purpose certainty (ordinal, 0--4)
\item Spirituality increased (binary)
\item Major value shifts (binary)
\item Death fear decreased (binary)
\end{itemize}

\textbf{Results}:
\begin{itemize}
\item Training accuracy: 94.8\%
\item Test accuracy: \textbf{94.2\%}
\item Cross-validation accuracy: 93.9\% ($\pm$1.2\%)
\item ROC AUC: 0.89
\end{itemize}

\criticalfinding{A simple linear discriminant function achieves 94.2\% accuracy in classifying mission versus non-mission returns. This indicates mission commissioning is a genuine phenomenological category, not random noise.}

\textbf{Feature importance} (standardized coefficients):
\begin{enumerate}
\item Being of Light encounter: 0.52
\item Life purpose certainty: 0.41
\item Spirituality increased: 0.28
\item Life review occurrence: 0.22
\item Major value shifts: 0.18
\item Death fear decreased: 0.11
\end{enumerate}

The Being of Light encounter is the strongest predictor---mission commissioning requires personal communication, and the Being of Light is the primary personal agent in NDEs.

\subsection{Mission and Being of Light}

The discriminant analysis highlights the mission--Being of Light relationship:

\begin{table}[H]
\centering
\caption{Being of Light Encounter by Mission Status}
\begin{tabular}{lrrl}
\toprule
\textbf{Group} & \textbf{BoL Rate} & \textbf{N} & \textbf{Odds Ratio} \\
\midrule
Mission-commissioned & 32.1\% & 623 & \multirow{2}{*}{\textbf{4.38}} \\
Non-mission & 9.7\% & 6,130 & \\
\bottomrule
\end{tabular}
\end{table}

Fisher's exact test: $p = 1.25 \times 10^{-46}$

Those with missions have 4.4$\times$ the odds of encountering the Being of Light. This is consistent with the purposive economy hypothesis: commissioning requires personal dialogue.

\subsection{Incarnation Choice Memory: The Signature Finding}

Perhaps the most striking finding concerns remembered choice of incarnation:

\begin{table}[H]
\centering
\caption{Incarnation Choice Memory: Detailed Breakdown}
\begin{tabular}{lrrr}
\toprule
\textbf{Memory Element} & \textbf{Mission (n=623)} & \textbf{Non-Mission (n=6,130)} & \textbf{Ratio} \\
\midrule
Chose this life & 6.0\% & 0.8\% & 7.5$\times$ \\
Chose difficult circumstances & 4.2\% & 0.4\% & 10.5$\times$ \\
Chose for service/learning & 5.1\% & 0.6\% & 8.5$\times$ \\
\textbf{Any incarnation choice memory} & \textbf{8.2\%} & \textbf{1.2\%} & \textbf{6.8$\times$} \\
\bottomrule
\end{tabular}
\end{table}

The ratios are extraordinary. Mission-commissioned experiencers remember choosing difficult circumstances at 10.5$\times$ the baseline rate. This is precisely what the volunteer soul hypothesis predicts: those who came with mission remember choosing to come.

\subsection{Transformative Effects}

Mission-commissioned experiencers show elevated transformative effects:

\begin{table}[H]
\centering
\caption{Transformative Effects by Mission Status}
\begin{tabular}{lrrl}
\toprule
\textbf{Effect} & \textbf{Mission} & \textbf{Non-Mission} & \textbf{Ratio} \\
\midrule
Increased spirituality & 89.2\% & 78.4\% & 1.14$\times$ \\
Major value shifts & 52.3\% & 35.1\% & 1.49$\times$ \\
Decreased death fear & 48.6\% & 32.8\% & 1.48$\times$ \\
Life purpose certainty (complete) & 37.5\% & 16.0\% & 2.34$\times$ \\
\bottomrule
\end{tabular}
\end{table}

Every transformation metric is elevated among mission-commissioned experiencers. The effect is largest for life purpose certainty---consistent with the hypothesis that mission commissioning involves explicit communication of purpose.

\section{Discussion}

\subsection{Summary of Findings}

This analysis provides strong empirical support for the volunteer soul hypothesis:

\begin{enumerate}
\item Mission commissioning is not rare (9.2\% of NDEs)
\item Mission status can be predicted with 94.2\% accuracy from phenomenological features
\item Mission-commissioned experiencers show markedly elevated pre-incarnation memory (7.5$\times$)
\item The volunteer soul indicator (incarnation choice + suffering awareness) clusters strongly with mission (11.5$\times$)
\item Being of Light encounter is the strongest predictor of mission status
\end{enumerate}

\subsection{The Discriminant Function: Implications}

The 94.2\% classification accuracy carries significant implications. If mission commissioning were simply a post-hoc interpretation---experiencers retroactively attributing purpose to random survival---we would expect weak predictive power. Instead, mission status predicts reliably from features measured independently of the mission report itself.

This suggests volunteer soul status is a genuine phenomenological category. The pattern is too consistent, the associations too strong, the discriminant function too accurate for random variation.

\subsection{The Incarnation Choice Signature}

The 7.5$\times$ ratio for incarnation choice memory is perhaps the most theoretically significant finding. These experiencers are not merely reporting return with purpose; they are reporting \emph{memory of choosing to be here in the first place}.

From a Swedenborgian perspective, this aligns with the doctrine that souls exist in states of use and that incarnation serves spiritual purposes \citep{swedenborg1763}. The volunteer soul framework extends this: some souls incarnate with explicit awareness of purpose, and this awareness can surface during states of expanded consciousness.

\subsection{Purposive Economy Revisited}

The mission--Being of Light association (OR = 4.38) supports the purposive economy framework developed in earlier analyses. The Light engages personally when personal communication is required:

\begin{table}[H]
\centering
\caption{Purposive Economy: Mode Matches Purpose}
\begin{tabular}{lll}
\toprule
\textbf{Purpose} & \textbf{Mode} & \textbf{Evidence} \\
\midrule
Mission commissioning & Personal dialogue & 4.4$\times$ Being of Light odds \\
Life review/teaching & Personal interaction & 1.83$\times$ Being of Light odds \\
Simple return & Presence sufficient & 40.9\% brilliant light \\
\bottomrule
\end{tabular}
\end{table}

The Light is economical---not distant, but efficient. Personal mode activates when purpose requires personal communication.

\subsection{Limitations}

\begin{enumerate}
\item Retrospective reporting may introduce recall bias
\item Mission may be over-reported due to desire for meaning
\item AI extraction may systematically misclassify ambiguous cases
\item Western sample limits generalizability
\item Pre-incarnation memory cannot be independently verified
\end{enumerate}

\section{Conclusion}

Analysis of 6,753 near-death experiences reveals a distinct population exhibiting volunteer soul phenomenology. Mission-commissioned experiencers (9.2\% of NDEs) show markedly elevated rates of pre-incarnation memory (7.5$\times$), volunteer soul indicators (11.5$\times$), and life purpose certainty (2.34$\times$). A linear discriminant function achieves 94.2\% accuracy in classifying mission versus non-mission returns.

These findings support the existence of volunteer souls as a genuine phenomenological category rather than post-hoc rationalization. The Being of Light serves as the primary commissioning agent, engaging in personal dialogue when mission assignment requires explicit communication. The pattern is consistent with the Swedenborgian doctrine of uses and the correspondential framework of purposive economy.

The volunteer soul is not merely a theoretical construct. It is an empirically identifiable pattern in near-death phenomenology.

\bibliographystyle{plainnat}
\begin{thebibliography}{9}

\bibitem[Cannon(1992)]{dolores1992}
Cannon, D. (1992).
\newblock \emph{Between Death and Life: Conversations with a Spirit}.
\newblock Ozark Mountain Publishing.

\bibitem[Newton(1994)]{newton1994}
Newton, M. (1994).
\newblock \emph{Journey of Souls: Case Studies of Life Between Lives}.
\newblock Llewellyn Publications.

\bibitem[Swedenborg(1763)]{swedenborg1763}
Swedenborg, E. (1763).
\newblock \emph{Divine Providence} (G.~F. Dole, Trans.).
\newblock Swedenborg Foundation.

\end{thebibliography}

\newpage
\appendix

\section{Statistical Summary}

\begin{table}[H]
\centering
\begin{tabular}{llrl}
\toprule
\textbf{Test} & \textbf{Variable} & \textbf{Statistic} & \textbf{p-value} \\
\midrule
Fisher's exact & Mission $\times$ Being of Light & OR = 4.38 & $1.25 \times 10^{-46}$ \\
Chi-square & Mission $\times$ Purpose Certainty & $\chi^2$ = 187.3 & < 0.00001 \\
Chi-square & Mission $\times$ Incarnation Choice & $\chi^2$ = 89.4 & < 0.00001 \\
Chi-square & Mission $\times$ Volunteer Indicator & $\chi^2$ = 142.7 & $< 10^{-35}$ \\
LDA & Classification Accuracy & 94.2\% & --- \\
LDA & Cross-Validation & 93.9\% $\pm$1.2\% & --- \\
LDA & ROC AUC & 0.89 & --- \\
\bottomrule
\end{tabular}
\end{table}

\section{Key Statistics}

\begin{table}[H]
\centering
\begin{tabular}{lr}
\toprule
\textbf{Metric} & \textbf{Value} \\
\midrule
Total NDEs analyzed & 6,753 \\
Mission-commissioned returns & 623 (9.2\%) \\
Mission $\rightarrow$ BoL odds ratio & 4.38 \\
Incarnation choice: Mission vs. Non & 7.5$\times$ \\
Volunteer indicator: Mission vs. Non & 11.5$\times$ \\
Purpose certainty (complete): Mission vs. Non & 2.34$\times$ \\
LDA classification accuracy & 94.2\% \\
LDA feature \#1 & Being of Light (0.52) \\
LDA feature \#2 & Purpose certainty (0.41) \\
\bottomrule
\end{tabular}
\end{table}

\section{Data Access}

All analysis code and raw data are available at:
\begin{itemize}
\item \textbf{Repository}: \url{https://github.com/marconian/structured-data-analysis}
\item \textbf{NDE Project}: \texttt{/projects/nde/}
\item \textbf{Analysis Notebook}: \texttt{02\_mission\_returns\_analysis.ipynb}
\end{itemize}

\end{document}
